%%%%%%%%%%%%%%%%%%%%%%%%%%%%%%%%%%%%%%%%%%%%%%%%%%%%%%%%%%%%%%%%%%%%%%%%%%%%%%%%
% BLOQUE 6 — ENTREGABLES ESPECÍFICOS
% Los 3 entregables específicos de la ficha TI-05, integrados en el
% cuerpo del informe (no como anexo). Todos en ESTE ÚNICO ARCHIVO (por
% preferencia del equipo, sin subcarpetas ni \input separados por
% entregable).
%%%%%%%%%%%%%%%%%%%%%%%%%%%%%%%%%%%%%%%%%%%%%%%%%%%%%%%%%%%%%%%%%%%%%%%%%%%%%%%%

\section{Entregables específicos}

% --- Entregable 1/3 (ficha TI-05) ---
\subsection{Modelo de costos: serverless vs. instancias siempre encendidas}

% TODO (equipo): modelo de costos que compare una arquitectura
% serverless con una de contenedores/instancias siempre encendidas,
% identificando el volumen de equilibrio en solicitudes por mes
% (fórmula, tabla o gráfico, según convenga).
%
% Recordatorio (punto 5 de las indicaciones sobre precios y fuentes):
% toda cifra debe indicar producto, moneda, región, fecha de consulta y
% si es precio de lista o cotización; verificar en la calculadora/
% página oficial del proveedor (no reutilizar cifras de terceros sin
% verificar). Si un proveedor no publica precios, declararlo
% expresamente ("solo por cotización").

% --- Entregable 2/3 (ficha TI-05) ---
\subsection{Prueba de concepto y mediciones de arranque en frío / latencia}

% TODO (equipo): elegir UNA de las dos opciones que exige la ficha:
%   (a) Prueba de concepto propia (código + mediciones), documentando
%       el entorno, el método de medición y los resultados; o
%   (b) Análisis de mediciones publicadas de arranque en frío/latencia,
%       citando la fuente y explicando su metodología.
% Si se incluye código propio, usar el entorno `minted` ya cargado por
% la clase (\begin{minted}{lenguaje} ... \end{minted}) y compilar con
% -shell-escape en Overleaf (Menu > Settings, o ya configurado según el
% motor del proyecto).

% --- Entregable 3/3 (ficha TI-05) ---
\subsection{Cuadro comparativo de plataformas}

% TODO (equipo): cuadro comparativo de AL MENOS 6 plataformas con sus
% límites técnicos (tiempo máximo de ejecución, tamaño máximo de
% payload, manejo de estado, portabilidad, etc.). Debe incluir
% plataformas de más de una categoría de la ficha (FaaS, contenedores
% sin servidor, cómputo en el borde, código abierto, datos sin
% servidor).
%
% OBLIGATORIO (punto 4 de las indicaciones): incorporar al menos DOS
% plataformas adicionales no mencionadas en la ficha (de ahí las 8
% filas de la tabla), con su justificación; o declarar expresamente que
% no se encontraron más alternativas relevantes y documentar la
% búsqueda (fuentes, criterios de descarte, fecha).
%
% Filas de la tabla: completar (mínimo 6 + 2 alternativas). Borrar las
% filas sobrantes o agregar más según corresponda.

\tabla{Cuadro comparativo de plataformas serverless y de borde}{cuadro-comparativo-plataformas}{|p{2.6cm}|p{2.3cm}|p{2cm}|p{2cm}|p{2.2cm}|p{3.4cm}|}{%
    \filaTabla{\textbf{Plataforma} & \textbf{Categoría} & \textbf{Tiempo máx. ejecución} & \textbf{Tamaño máx. payload} & \textbf{Manejo de estado} & \textbf{Otros límites / notas}}
    \filaTabla{ & & & & & }
    \filaTabla{ & & & & & }
    \filaTabla{ & & & & & }
    \filaTabla{ & & & & & }
    \filaTabla{ & & & & & }
    \filaTabla{ & & & & & }
    \filaTabla{ & & & & & }
    \filaTabla{ & & & & & }
}

\clearpage
